\subsection{Hartogs' ordinal}
The Hartogs' ordinal is a special ordinal that can be used to prove the existence of a 
well ordering for any set under certain assumptions. It can be defined using axioms of ZF alone. 
Specifically, its existence does not depend on the Axiom of Choice (AC). \\

\begin{definition}[Hartogs' ordinal]
    The Hartogs' ordinal $\alpha$ of a set $A$ is the smallest ordinal $\alpha$ such that there 
    exists an injection from $A$ to $\alpha$. In other words, it is the smallest ordinal that is 
    "larger" than $A$.
\end{definition}

\begin{theorem}[Hartogs' theorem]
    For every set $A$ there exists an Hartogs' ordinal.
\end{theorem}
\begin{proof}
    If $A$ is empty, then the Hartogs' ordinal is $0$. Otherwise, we can construct the Hartogs'
    ordinal by taking the set of all ordinals that are injectable into $A$. The set that is 
    constructed this way is itself an ordinal, and since it doesn't belong to the set itself 
    it must not be injectable into $A$. Therefore, it is the Hartogs' ordinal of $A$.
\end{proof}

\newpage